\documentclass[11pt,a4paper]{article}
\usepackage[margin=3cm]{geometry}
\usepackage{amsmath,amssymb,amsthm,mathrsfs}
\usepackage[T1]{fontenc}
\usepackage{microtype}
\usepackage{booktabs}
\usepackage[colorlinks=true,linkcolor=blue,citecolor=blue]{hyperref}

\theoremstyle{plain}
\newtheorem{theorem}{Theorem}[section]
\newtheorem{proposition}[theorem]{Proposition}
\newtheorem{lemma}[theorem]{Lemma}
\newtheorem{corollary}[theorem]{Corollary}
\theoremstyle{definition}
\newtheorem{remark}[theorem]{Remark}

\newcommand{\A}{\mathfrak{A}}
\newcommand{\HH}{\mathcal{H}}
\newcommand{\G}{\mathcal{G}}
\newcommand{\R}{\mathbb{R}}
\newcommand{\Om}{\Omega}
\newcommand{\ip}[2]{\langle #1,#2\rangle}
\DeclareMathOperator{\spec}{spec}
\DeclareMathOperator{\dist}{dist}
\DeclareMathOperator{\supp}{supp}

\title{M2: the falsifiable checkpoint\\[2pt]
\large Which mass sets the rate of $\omega_g\to\omega_\infty$?}
\author{}
\date{}

\begin{document}
\maketitle

\begin{abstract}
We run the M0 architecture on the exactly solvable quadratic model $P(\xi)=\xi^2$, where Rosen
\cite{Rosen72} identifies the limit, and we extract the convergence \emph{rate}, which Rosen does
not. Two findings.

\smallskip
\noindent\textbf{(1) Consistency: pass.} For the quadratic model $(\A,\alpha)$ \emph{is} the free
field of mass $m_1=(m_0^2+\lambda)^{1/2}$, whose locally normal ground state is unique. Hence M0
plus uniqueness reproduces Rosen's Theorem~5.2 --- $\omega_g\to\omega_{m_1}$ in norm on every
local algebra --- with no Bogoliubov transformation, no dressing operator and no Shale condition.

\smallskip
\noindent\textbf{(2) Rate: the strategy's conjecture was wrong by a factor $2$.} The exponent is
governed by $2m_1$, not by $m_0$ and not by $m_1$:
\[
  \bigl\|(\omega_g-\omega_\infty)|_{\A(O)}\bigr\|\ \sim\ \mathrm{poly}(d)\,e^{-2m_1 d},
  \qquad d=\dist(O,\{g<1\}).
\]
The mass is the \emph{physical} one because the resolvent kernels decay at the local Agmon rate,
which is $m_1$ inside $\{g=1\}$; the factor $2$ is a \emph{round trip} $O\to\supp(1-g)\to O$
forced by the exact resolvent identity. Confirmed numerically to $\sim1\%$ over six decades
(\S\ref{sec:num}).
\end{abstract}

\section{The model and the two questions}

Take $P(\xi)=\xi^2$, i.e.
\[
  H(g)=H_0+\tfrac{\lambda}{2}\int_\R g(x):\!\varphi_0(x)^2\!:dx-E(g),\qquad \lambda>0,
\]
and set
\[
  A:=-\Delta+m_0^2+\lambda g,\qquad B:=-\Delta+m_1^2,\qquad m_1^2:=m_0^2+\lambda ,
\]
so that $\mu_g:=A^{1/2}$ and $\mu_1:=B^{1/2}$ are the one-particle energies. Both $\omega_g$ and
the candidate limit $\omega_\infty$ are quasi-free:
\begin{equation}\label{eq:qf}
  \omega_g\bigl(W(f,h)\bigr)=\exp\Bigl(-\tfrac14\bigl[\ip{f}{\mu_g^{-1}f}+\ip{h}{\mu_g h}\bigr]\Bigr),
\end{equation}
and likewise for $\omega_\infty$ with $\mu_1$. Everything therefore reduces to the two operator
differences
\begin{equation}\label{eq:Deltas}
  \Delta_\varphi:=\mu_g^{-1}-\mu_1^{-1},\qquad
  \Delta_\pi:=\mu_g-\mu_1 .
\end{equation}

Fix a bounded open $O$ and let $I$ be an interval with $g\equiv1$ on $I$ and $O\subset I$; put
$d:=\dist(O,\R\setminus I)$.

\section{Consistency: the architecture reproves Rosen}\label{sec:consistency}

\begin{lemma}\label{lem:freeeq}
For the quadratic model, the Glimm--Jaffe automorphism group $\alpha$ of $\A$ is the free
time evolution of mass $m_1$.
\end{lemma}

\begin{proof}
Let $A\in\A(O)$ and $t\in\R$, and pick $g\in\G$ with $g\equiv1$ on $O_{|t|}$. By
Fact~1.3(ii) of \cite{M0}, $\alpha_t(A)=e^{itH(g)}Ae^{-itH(g)}$. The Heisenberg equations for
$H(g)$ are the linear field equations $\partial_t^2\varphi=\Delta\varphi-(m_0^2+\lambda g)\varphi$,
whose solution on $O_{|t|}$ depends, by finite propagation speed, only on the coefficients in
$O_{|t|}$, where $m_0^2+\lambda g\equiv m_1^2$. Hence $\alpha_t$ restricted to $\A(O)$ is the free
mass-$m_1$ evolution. As $O$ and $t$ were arbitrary and $\bigcup_O\A(O)$ is norm-dense, the claim
follows.
\end{proof}

\begin{lemma}\label{lem:uniquegs}
The free field of mass $m_1>0$ has a unique locally normal ground state, namely the quasi-free
state $\omega_{m_1}$ with covariances $(2\mu_1)^{-1}$ and $\mu_1/2$.
\end{lemma}

\begin{proof}[Proof sketch]
Let $\omega$ be a ground state. In the GNS representation the generator $H_\omega\ge0$ annihilates
$\Om_\omega$, and the free dynamics acts on the one-particle test-function space by the symplectic
group generated by $\mu_1$. Positivity of $H_\omega$ forces the two-point function to be the
$\mu_1$-vacuum one, exactly as in the classical argument for the free field; local normality
excludes the non-regular ground states carried by non-normal representations. Uniqueness of the
one-particle covariance plus quasi-freeness of $\omega_{m_1}$ then gives $\omega=\omega_{m_1}$.
\end{proof}

\begin{corollary}[consistency check: \textbf{pass}]\label{cor:rosen}
Assume the local number bound \textup{(H2)} of \cite{M0}. Then, in the quadratic model,
$\omega_g\to\omega_{m_1}$ in norm on every $\A(O)$.
\end{corollary}

\begin{proof}
Combine \cite[Thm.~4.2, Cor.~4.4]{M0} with Lemmas~\ref{lem:freeeq} and~\ref{lem:uniquegs}.
\end{proof}

This is exactly Rosen's Theorem~5.2 \cite[\S5]{Rosen72}. The point is the route: Rosen identifies
the limit through the explicit dressing transformation $T_g=e^{iL}\Gamma(u)$ and the Shale
condition $\ln p\in\mathrm{HS}$, none of which appears above. The M0 architecture obtains the same
conclusion from the exact ground-state identity plus uniqueness, i.e.\ from precisely the two
ingredients that are also available --- in principle --- for $\varphi^4_2$.

\section{The rate}\label{sec:rate}

\subsection{The exact resolvent identity, and why a round trip appears}

Since $A-B=-\lambda(1-g)$ and $(A+s)^{-1}-(B+s)^{-1}=-(A+s)^{-1}(A-B)(B+s)^{-1}$, the standard
representations $S^{-1/2}=\frac1\pi\int_0^\infty s^{-1/2}(S+s)^{-1}ds$ and
$S^{1/2}=\frac1\pi\int_0^\infty s^{-1/2}S(S+s)^{-1}ds$ give the \emph{exact} identities
\begin{align}
  \Delta_\varphi &= \frac{\lambda}{\pi}\int_0^\infty s^{-1/2}\,(A+s)^{-1}\,(1-g)\,(B+s)^{-1}\,ds\ \ge 0,
  \label{eq:idphi}\\[2pt]
  \Delta_\pi &= -\frac{\lambda}{\pi}\int_0^\infty s^{+1/2}\,(B+s)^{-1}\,(1-g)\,(A+s)^{-1}\,ds\ \le 0 .
  \label{eq:idpi}
\end{align}
These are not perturbative expansions: they are identities. The signs agree with
$\mu_g\le\mu_1$ and are confirmed numerically in \S\ref{sec:num}.

The structural consequence is immediate and is the whole content of M2. The factor $1-g$ is
supported in $\R\setminus I$, at distance $\ge d$ from $O$. Hence the kernel of
$P_O\Delta_\varphi P_O$ is a sum over paths that leave $O$, reach $\supp(1-g)$, and return: the
exponential cost is paid \emph{twice}.

\subsection{Which mass}

\begin{lemma}[local Agmon decay]\label{lem:agmon}
For $s\ge0$ let $\kappa_j(s):=(m_j^2+s)^{1/2}$. Then
\[
  (B+s)^{-1}(x,y)=\frac{1}{2\kappa_1}e^{-\kappa_1|x-y|}\quad\text{(exactly, in one dimension),}
\]
and there is $C$ independent of $s,g$ with
\[
  \bigl|(A+s)^{-1}(x,z)\bigr|\ \le\ \frac{C}{\kappa_0}\,
  \exp\Bigl(-\int_x^z\bigl(m_0^2+\lambda g(u)+s\bigr)^{1/2}du\Bigr).
\]
\end{lemma}

The second bound is the standard Agmon/Combes--Thomas estimate; the point is that the decay rate
is the \emph{local} one. For $x\in O$ and $z\notin I$ the path crosses $I$, where
$m_0^2+\lambda g\equiv m_1^2$, so
\[
  \int_x^z\bigl(m_0^2+\lambda g+s\bigr)^{1/2}\,du\ \ge\ \kappa_1(s)\,d+\kappa_0(s)\bigl(\dist(z,I)\bigr).
\]

\begin{proposition}\label{prop:rate}
There is $C=C(m_0,\lambda,|O|)$ such that for all $d\ge1$,
\[
  \bigl\|P_O\Delta_\varphi P_O\bigr\|\le C\,d^{-1/2}e^{-2m_1d},
  \qquad
  \bigl\|P_O\Delta_\pi P_O\bigr\|\le C\,d^{-3/2}e^{-2m_1d} .
\]
Moreover the exponent is \emph{sharp}: by \eqref{eq:qf},
$\omega_g(\varphi_0(f)^2)-\omega_\infty(\varphi_0(f)^2)=\tfrac12\ip{f}{\Delta_\varphi f}$ for
$\supp f\subset O$, so $\|(\omega_g-\omega_\infty)|_{\A(O)}\|$ is bounded below by a constant
times $\|P_O\Delta_\varphi P_O\|$ as well as above.
\end{proposition}

\begin{proof}[Proof sketch]
Insert Lemma~\ref{lem:agmon} into \eqref{eq:idphi}. The $z$-integral over $\R\setminus I$ gives a
factor $\le 2e^{-\kappa_1 d}/(\kappa_0+\kappa_1)$, so the two exponentials combine to
$e^{-2\kappa_1(s)d}$. Substituting $\kappa=\kappa_1(s)$, $s=\kappa^2-m_1^2$, $ds=2\kappa\,d\kappa$,
\[
  \int_0^\infty s^{-1/2}e^{-2\kappa_1(s)d}\,\frac{ds}{\kappa_1^2(\kappa_0+\kappa_1)}
  \;\asymp\;\int_{m_1}^\infty\frac{2\kappa\,e^{-2\kappa d}}{(\kappa^2-m_1^2)^{1/2}\kappa^2(\kappa_0+\kappa)}\,d\kappa
  \;\asymp\;C\,d^{-1/2}e^{-2m_1d},
\]
the last step by Watson's lemma at the endpoint $\kappa=m_1$, where the integrand has an inverse
square-root singularity. For \eqref{eq:idpi} the weight $s^{+1/2}$ replaces $s^{-1/2}$, turning
$(\kappa^2-m_1^2)^{-1/2}$ into $(\kappa^2-m_1^2)^{+1/2}$ and hence $d^{-1/2}$ into $d^{-3/2}$.
\end{proof}

\begin{remark}[what M2 falsifies]
The strategy note conjectured $\|(\omega_g-\omega_\infty)|_{\A(O)}\|\lesssim e^{-m\,d}$ with $m$
the mass gap. Proposition~\ref{prop:rate} says the correct statement carries a factor $2$ in the
exponent. Both features generalise beyond the quadratic model: the round trip is forced by the
resolvent identity, i.e.\ by the fact that the perturbation is \emph{not} between $O$ and the
observable but between $O$ and the region where the two Hamiltonians differ; and the relevant
mass is the local one inside $\{g=1\}$, which is the physical mass, not the bare mass $m_0$.
Any quasi-adiabatic argument at Milestone M4 must reproduce $2m$, not $m$.
\end{remark}

\section{Numerical confirmation}\label{sec:num}

Script: \texttt{m2\_rate.py}; output: \texttt{m2\_rate.out}. Parameters $m_0=1$, $\lambda=3$, hence
$m_1=2$ and the four candidate exponents are $m_0=1$, $2m_0=2$, $m_1=2$, $2m_1=4$. Sharp cutoff
$g=\mathbf 1_{[-\ell,\ell]}$, box $[-25,25]$, spacing $h=0.025$ (and $h=0.0166$ for a refinement
check), $O=\{0\}$ so $d=\ell$. Kernels evaluated at the origin; one eigendecomposition per $\ell$.

\begin{center}
\begin{tabular}{lrrr}
\toprule
 & fitted exponent $k$ & prefactor power $p$ & max.\ residual\\
\midrule
$\Delta_\varphi$, joint fit $\log|y|=c-kd-p\log d$ & $4.0343$ & $0.288$ & $1.1\cdot10^{-3}$\\
$\Delta_\pi$,\ \ joint fit                          & $4.0767$ & $1.019$ & $4.0\cdot10^{-3}$\\
\midrule
prediction & $2m_1=4$ & $1/2$, $3/2$ & ---\\
\bottomrule
\end{tabular}
\end{center}

Prefactor-free evidence: the local slopes $k(d)=-d\log|y|/dd$, which for
$|y|\sim d^{-p}e^{-kd}$ satisfy $k(d)=k+p/d$, decrease monotonically over $d\in[1,4]$ from
$4.285$ to $4.105$ ($\Delta_\varphi$) and from $4.965$ to $4.322$ ($\Delta_\pi$); subtracting
$p/d$ with the fitted $p$ gives $4.031$ and $4.060$. Refining $h$ by $1.5\times$ moves the
$\Delta_\varphi$ exponent by $<0.03$. Both sign predictions of \eqref{eq:idphi}--\eqref{eq:idpi}
hold at every data point.

\medskip
\noindent\textbf{Verdict.} $k=4.03\pm0.05$ against candidates $\{1,2,2,4\}$: the exponent is
$2m_1$, decisively. The polynomial prefactor powers come out $0.29$ and $1.02$ against the
predicted $0.5$ and $1.5$; over $d\in[1,4]$ the asymptotic expansion has relative corrections of
order $1/d$, so these are consistent in sign and magnitude but are \emph{not} sharply confirmed.
The exponent, which is what M2 was designed to test, does not depend on them.

\section{Consequences for the programme}

\begin{enumerate}
\item \textbf{The architecture is sound.} Corollary~\ref{cor:rosen} recovers a known theorem from
      M0 plus uniqueness, using no structure special to the quadratic case beyond the uniqueness
      input itself. This is the proof-of-concept the strategy needed.
\item \textbf{Correct the rate in the strategy}: $e^{-m\,d}\rightsquigarrow e^{-2m\,d}$
      throughout \S6, and correspondingly in the target statement of Theorem~A.
\item \textbf{Rosen's estimates are not sharp for this question.} His Lemma~5.1 carries
      $e^{-m_0|j|}$ because it bounds the \emph{free, mass-$m_0$} kernel; it is used only to prove
      $\|(\mu_g-\mu_1)\xi\|\to0$ qualitatively. The sharp local rate involves $m_1$ and the factor
      $2$. This is consistent with Rosen never stating a rate.
\item \textbf{A target for M4.} Any quasi-adiabatic/LPPL argument must produce $2m$; if it
      produces $m$, it is lossy by a factor $2$ in the exponent, and if it produces $m_0$ it is
      wrong.
\end{enumerate}

\begin{thebibliography}{9}
\bibitem{M0} \emph{Limit points of the finite-volume ground states of $P(\varphi)_2$ are ground
states of the quasi-local dynamics}, Milestone M0, \texttt{../m0/m0.tex}.
\bibitem{Rosen72} L.~Rosen, \emph{Renormalization of the Hilbert space in the mass shift model},
J.\ Math.\ Phys.\ \textbf{13} (1972) 918--927.
\end{thebibliography}

\end{document}
